% STATUS: written 2026-10-21, recalled results as is (one environment each, source kind and whole statement)
% Sources read in full: Fosner and Ilisevic 2008 (pp. 415--427, Sections 1 and 5), Vukman 1997 (pp. 1--6, read from OCR and checked against the page images).
\section{Jordan *-derivations and Jordan left derivations}\label{sec:jordan-type:star-left}

Section~\ref{sec:jordan-type:triple-and-generalized} ended with the triple form of the product rule.
Two further changes of the rule remain. Fošner and Ilišević use an involution
\cite[p.~416]{FosnerIlisevic2008}. Vukman studies Jordan left derivations
\cite[p.~1]{Vukman1997}.

Fošner and Ilišević start from a $\ast$-ring, that is, a ring $R$ with an involution $\ast$
\cite[Section~1, p.~415]{FosnerIlisevic2008}. An involution is an additive mapping with
$(a^{\ast})^{\ast}=a$ and $(ab)^{\ast}=b^{\ast}a^{\ast}$ for all $a,b\in R$
\cite[p.~415]{FosnerIlisevic2008}.

\begin{definition}[{\cite[Section~1, p.~416]{FosnerIlisevic2008}}]\label{def:jordan:star}
Let $R$ be a $\ast$-ring and let $M$ be an $R$-bimodule. An additive mapping $\delta\colon R\to M$ is
called a Jordan $\ast$-derivation if $\delta(x^{2})=\delta(x)x^{\ast}+x\delta(x)$ holds for all
$x\in R$. A Jordan triple $\ast$-derivation is an additive mapping $\delta\colon R\to M$ with the
property $\delta(xyx)=\delta(x)y^{\ast}x^{\ast}+x\delta(y)x^{\ast}+xy\delta(x)$ for all $x,y\in R$.
\end{definition}

On a $2$-torsion free $\ast$-ring every Jordan $\ast$-derivation is a Jordan triple $\ast$-derivation
\cite[p.~417]{FosnerIlisevic2008}. For real Banach $\ast$-algebras with identity the converse holds
\cite[p.~417]{FosnerIlisevic2008}. A recent paper proved the converse for $6$-torsion free semiprime
$\ast$-rings \cite[p.~417]{FosnerIlisevic2008}. Fošner and Ilišević extend this result to $2$-torsion free semiprime
$\ast$-rings \cite[p.~417]{FosnerIlisevic2008}. They rest on a lemma whose proof is found in an earlier paper \cite[p.~425]{FosnerIlisevic2008}.

\begin{lemma}[{\cite[Lemma~5.1, p.~425]{FosnerIlisevic2008}}]\label{lem:jordan:star-semiprime}
If $R$ is a semiprime $\ast$-ring, then $yxy^{\ast}=0$ for all $y\in R$ (where $x\in R$) implies
$x=0$.
\end{lemma}

The lemma gives the main result of the section.

\begin{theorem}[{\cite[Theorem~5.2, p.~425]{FosnerIlisevic2008}}]\label{thm:jordan:star-triple}
Let $R$ be a $2$-torsion free semiprime $\ast$-ring and let $\delta\colon R\to R$ be an additive
mapping. Then the following conditions are mutually equivalent:
(a) $\delta$ is a Jordan $\ast$-derivation;
(b) $\delta(xyx)=\delta(x)y^{\ast}x^{\ast}+x\delta(y)x^{\ast}+xy\delta(x)$ $(x,y\in R)$.
\end{theorem}

The generalized version uses a second map $E$ that is related to $\delta$
\cite[p.~417]{FosnerIlisevic2008}. An additive mapping $E$ is called a generalized Jordan
$\ast$-derivation if a Jordan $\ast$-derivation $\delta$ exists with $E(x^{2})=E(x)x^{\ast}+x\delta(x)$
\cite[p.~417]{FosnerIlisevic2008}.

\begin{theorem}[{\cite[Theorem~5.4, p.~426]{FosnerIlisevic2008}}]\label{thm:jordan:star-generalized}
Let $R$ be a $2$-torsion free semiprime $\ast$-ring and let $E,\delta\colon R\to R$ be additive
mappings. Then the following conditions are mutually equivalent:
(a) for all $x,y\in R$, $E(xyx)=E(x)y^{\ast}x^{\ast}+x\delta(y)x^{\ast}+xy\delta(x)$ and
$\delta(xyx)=\delta(x)y^{\ast}x^{\ast}+x\delta(y)x^{\ast}+xy\delta(x)$;
(b) for all $x\in R$, $E(x^{2})=E(x)x^{\ast}+x\delta(x)$ and
$\delta(x^{2})=\delta(x)x^{\ast}+x\delta(x)$.
\end{theorem}

Let $R$ be an associative ring and $D\colon R\to R$ an
additive mapping \cite[p.~1]{Vukman1997}.

\begin{definition}[{\cite[p.~1]{Vukman1997}}]\label{def:jordan:left}
$D$ is called a derivation (respectively a left derivation) if $D(xy)=D(x)y+xD(y)$ (respectively
$D(xy)=xD(y)+yD(x)$) holds for all $x,y\in R$. And $D$ is called a Jordan derivation (respectively a left
Jordan derivation) in case $D(x^{2})=D(x)x+xD(x)$ (respectively $D(x^{2})=2xD(x)$) for all $x\in R$.
\end{definition}

Vukman notes that a left derivation of a noncommutative prime ring is zero
\cite[p.~1]{Vukman1997}. He also records a theorem of Brešar and himself
\cite[p.~1]{Vukman1997}. It says that a nonzero Jordan left derivation on a prime ring of
characteristic different from $2$ and $3$ forces the ring to be commutative
\cite[p.~1]{Vukman1997}. Vukman states three lemmas for his theorems \cite[p.~2]{Vukman1997}. The
first holds in any $2$-torsion free ring.

\begin{lemma}[{\cite[Lemma~3, p.~2]{Vukman1997}}]\label{lem:jordan:left-identities}
Let $R$ be a $2$-torsion free ring. If $D\colon R\to R$ is a Jordan left derivation then for all
$x,y\in R$: (a) $D(xy+yx)=2xD(y)+2yD(x)$, (b) $D(xyx)=x^{2}D(y)+3xyD(x)-yxD(x)$.
\end{lemma}

The second lemma treats prime rings \cite[p.~2]{Vukman1997}.

\begin{lemma}[{\cite[Lemma~5, p.~2]{Vukman1997}}]\label{lem:jordan:left-prime}
Let $R$ a noncommutative prime ring of characteristic different from two, and let $D\colon R\to R$ be a
left Jordan derivation. In this case $D=0$.
\end{lemma}

Vukman combines these with a result of Giambruno and Herstein on nilpotent derivations
\cite[p.~1]{Vukman1997}. His first theorem treats semiprime rings.

\begin{theorem}[{\cite[Theorem~1, p.~1]{Vukman1997}}]\label{thm:jordan:left-nilpotent}
Let $R$ be a $2$-torsion free semiprime ring and $D\colon R\to R$ a Jordan left derivation. If there
exists a positive integer $n$ such that $D(x)^{n}=0$ for all $x\in R$, then $D=0$.
\end{theorem}

Deng treats the left Jordan condition on its own, without the prime-ring setting
Vukman uses, and the map it governs is the one the definition above names
\cite[Abstract]{Deng1992JordanLeft}.

The Jordan left condition can be imposed at a product rather than at every
element. On a Banach algebra a continuous linear map satisfying
$2a\delta(b)+2b\delta(a)=\delta(z)$ whenever $ab+ba=z$ is characterized for
$z=1$, and then for $z$ a nontrivial idempotent of a reflexive algebra
$\operatorname{Alg}L$ \cite[Abstract]{ContinuousLinearReflexive2013}.

Strong commutativity preservation has a skew form. For two generalized skew
derivations $F$ and $G$ of a prime ring with the same associated automorphism
$\alpha$, if $[F(x),\alpha(y)]=G([x,y])$ on the values of a non-central
polynomial, then $F$ is a scalar multiple of $\alpha$ plus a suitable map
\cite[Abstract]{AutomorphismsGeneralizedSkew2016}.

